% \chapter{Validation}
% \label{chap:validation}

% The purpose of this chapter is to evaluate whether the proposed solution achieves the objectives defined in \Cref{chap:proposed-solution}. The evaluation combines objective measurements of system behaviour with subjective feedback collected from users. Together, these perspectives provide a comprehensive assessment of the prototype's effectiveness as an \textit{IEC 61499} development environment.

% The evaluation is organised into five complementary areas. Functional validation verifies that the prototype correctly supports the modelling, deployment, and monitoring of \textit{IEC 61499} applications. Performance evaluation analyses workflow efficiency and compares key development tasks against existing tools. Collaboration evaluation examines the benefits of the real-time collaborative architecture and contrasts it with traditional version-control-based workflows. Usability evaluation focuses on user perception and experience, including a study of existing \textit{IEC 61499} tooling. Finally, the results are discussed in the context of the research objectives and identified limitations.

\section{Evaluation Goals}

The evaluation aims to answer the following questions:

\begin{enumerate}
\item Can the proposed solution successfully support the complete development workflow of an \textit{IEC 61499} application?
\item Does the proposed architecture provide advantages over existing desktop-based development environments?
\item Does the integration of real-time collaboration improve collaborative development workflows?
\item Do users perceive the proposed solution as accessible and easy to use?
\end{enumerate}

To answer these questions, both quantitative and qualitative evaluation methods were employed.

\section{Experimental Setup}

% A common test application was used throughout the validation process. The same application was implemented using both the developed prototype and \textit{4diac IDE}, allowing direct comparison between environments.

% \subsection{Available Resources}

% The following resources were made available to the tested \acp{IDE}:

% \begin{itemize}
% \item One \textit{DINASORE} runtime environment.
% \item Three function block definitions obtained from the official \textit{DINASORE} tutorial \textit{Sensorization "Hello World!"}:
% \begin{itemize}
% \item \texttt{SENSOR_SIMULATOR.fbt} and \texttt{SENSOR_SIMULATOR.py}
% \item \texttt{MOVING_AVERAGE.fbt} and \texttt{MOVING_AVERAGE.py}
% \item \texttt{CONCATENATE_REALS.fbt} and \texttt{CONCATENATE_REALS.py}
% \end{itemize}
% \end{itemize}

% The selected function blocks provide a representative example of a small \textit{IEC 61499} application while remaining sufficiently simple for controlled testing.

% \subsection{Test Application}

% The application shown in \Cref{fig:test-program} was implemented and deployed during the validation process.

% \begin{figure}[H]
% \centering
% \includegraphics[width=0.8\linewidth]{test-program.jpg}
% \caption{Test application used during validation}
% \label{fig:test-program}
% \end{figure}

% The application consists of a sensor simulation function block whose output is processed by a moving average filter before being formatted for display.

% \subsection{Expected Behaviour}

% After deployment, the application should execute correctly on the target \textit{DINASORE} instance and continuously produce output values. These values should be observable through the monitoring facilities provided by the \ac{IDE}.

% \section{Evaluation Metrics}

% The evaluation combines both objective and subjective metrics.

\subsection{Functional Metrics}

Functional validation considers:

\begin{itemize}
\item Successful import of function block definitions
\item Correct graph modelling behaviour
\item Successful deployment to \textit{DINASORE}
\item Correct creation of deployment commands
\item Successful monitoring of runtime values
\end{itemize}

\subsection{Performance Metrics}

Performance evaluation considers:

\begin{itemize}
\item Installation complexity
\item Startup time
\item Number of manual configuration steps
\item Deployment workflow complexity
\end{itemize}

\subsection{Collaboration Metrics}

Collaboration capabilities are evaluated through:

\begin{itemize}
\item Support for concurrent editing
\item Conflict handling mechanisms
\item Need for manual synchronisation
\end{itemize}

\subsection{Usability Metrics}

Usability is assessed through user questionnaires based on:

\begin{itemize}
\item Perceived task difficulty
\item Interface intuitiveness
\item User satisfaction
\item Open-ended qualitative feedback
\end{itemize}
